\documentclass[10pt]{article}
\usepackage[margin=1in]{geometry}
\usepackage{amsmath,amssymb,mathtools}
\usepackage{booktabs}
\usepackage{xcolor}
\usepackage[hidelinks]{hyperref}
\usepackage{microtype}
\usepackage{fancyhdr}
\definecolor{navy}{RGB}{24,55,95}
\pagestyle{fancy}
\setlength{\headheight}{14pt}
\fancyhf{}
\lhead{Symmetric groups}
\rhead{Exact finite-$n$ proof and verification}
\cfoot{\thepage}
\newcommand{\Sym}{\operatorname{Sym}}
\newcommand{\M}{\mathcal M}
\title{\color{navy}\textbf{The Exact Finite-$n$ $\sigma$-MGF for $S_n$}\\[2mm]
\large Proof, orbit interpretation, and matrix verification}
\author{Computational proof companion}
\date{17 July 2026}

\begin{document}
\maketitle
\begin{abstract}
For the permutation representation $V=\mathbb C^n$, this note proves the exact
cycle-type formula for $\M_{S_n}$, identifies its coefficients with vector
partitions, and explains the stable range.  The accompanying marimo notebook
constructs all permutation matrices and compares a full Reynolds-projector
rank with both formulas.  All representative tests passed.
\end{abstract}

\section{Molien coefficient as an invariant dimension}
For $\tau=(\tau_1,\ldots,\tau_s)$ let
$W_\tau=\bigotimes_b\Sym^{\tau_b}(\mathbb C^n)$.  Since
\[
 \det(I-tg)^{-1}=\sum_{q\ge0}\chi_{\Sym^qV}(g)t^q,
\]
the coefficient of $m_\tau$ in the generalized Molien series is
\[
 [m_\tau]\M_{S_n}=\frac1{n!}\sum_{g\in S_n}\chi_{W_\tau}(g)
 =\operatorname{tr}\left(\frac1{n!}\sum_{g\in S_n}\rho_{W_\tau}(g)\right)
 =\dim W_\tau^{S_n}.
\]
The middle operator is the Reynolds projector.

\section{Exact cycle-type formula}
Write a cycle type as $\lambda=(1^{m_1}2^{m_2}\cdots)\vdash n$ and set
\[
 z_\lambda=\prod_{d\ge1}d^{m_d}m_d!.
\]
A $d$-cycle has all $d$th roots of unity as eigenvalues, so
$\det(I-tP_d)=1-t^d$.  A permutation of type $\lambda$ therefore contributes
\[
 \prod_i\prod_{d\ge1}(1-t_i^d)^{-m_d(\lambda)}.
\]
The conjugacy class has size $n!/z_\lambda$, proving
\[
 \boxed{\M_{S_n}=\sum_{\lambda\vdash n}\frac1{z_\lambda}
 \prod_i\prod_{d\ge1}(1-t_i^d)^{-m_d(\lambda)}.}
\]
This statement is exact for every $n$ and in every degree; it is not a stable
approximation.

\section{Exact coefficient formula}
A monomial-basis vector of $W_\tau$ is encoded by a labeling
\[
 \ell:\{1,\ldots,n\}\longrightarrow\mathbb Z_{\ge0}^s,
 \qquad \sum_{j=1}^n\ell(j)=\tau
\]
coordinatewise.  Two labelings lie in the same $S_n$-orbit exactly when the
multiplicities
$a_{\mathbf v}=\#\{j:\ell(j)=\mathbf v\}$ agree.  Since the dimension of the
invariant space of a permutation module equals its number of orbits,
\[
 \boxed{[m_\tau]\M_{S_n}=\#\left\{(a_{\mathbf v})_{\mathbf v\ge0}:
 a_{\mathbf v}\ge0,\quad \sum_{\mathbf v}a_{\mathbf v}=n,\quad
 \sum_{\mathbf v}a_{\mathbf v}\mathbf v=\tau\right\}.}
\]
Equivalently, this is the number of vector partitions of $\tau$ into at most
$n$ nonzero parts; zero parts fill the remaining slots.

Marking the number of labels with $u$ gives
\[
 \boxed{\sum_{n\ge0}u^n\M_{S_n}
 =\prod_{\mathbf v\in\mathbb Z_{\ge0}^{(\mathbb N)}}(1-ut^{\mathbf v})^{-1}.}
\]
The exponential formula for permutation cycles gives the equivalent form
\[
 \sum_{n\ge0}u^n\M_{S_n}
 =\exp\left(\sum_{d\ge1}\frac{u^d}{d}
 \prod_i\frac1{1-t_i^d}\right).
\]

\section{Stable range}
Separate the zero label.  Every nonzero vector has total degree at least one,
so a vector partition of $\tau$ uses at most $|\tau|$ nonzero parts.  If
$n\ge|\tau|$, zeros can always be added and the coefficient stabilizes:
\[
 [m_\tau]\M_{S_n}=[m_\tau]\prod_{\mathbf v\ne0}(1-t^{\mathbf v})^{-1}.
\]
When $n<|\tau|$, the exact finite-$n$ formula simply removes decompositions
requiring too many nonzero parts.

\section{Verification method and evidence}
For each $g\in S_n$, the notebook creates its permutation matrix, compresses
$g^{\otimes q}$ to an orthonormal occupancy basis of $\Sym^q(V)$, tensors the
selected blocks, and averages the resulting matrices.  The projector rank is
compared independently with the symmetric-power character average, cycle-type
coefficient extraction, and a dynamic-programming vector-partition count.

\begin{center}
\begin{tabular}{@{}lrrrrr@{}}
\toprule
Case & $\dim W_\tau$ & rank & character & cycle & orbit\\
\midrule
$S_2,\tau=(3)$ & 4 & 2 & 2 & 2 & 2\\
$S_3,\tau=(2,1)$ & 18 & 4 & 4 & 4 & 4\\
$S_4,\tau=(2,2)$ & 100 & 9 & 9 & 9 & 9\\
$S_5,\tau=(2,1)$ & 75 & 4 & 4 & 4 & 4\\
\bottomrule
\end{tabular}
\end{center}
The largest projector here is $100\times100$ and averages 24 explicitly
constructed matrices.  Idempotence residuals range from
$2.2\times10^{-16}$ to $1.1\times10^{-15}$.  These cases cover degree above
the stable threshold, multiblock tensor products, and two different values of
$n$ for the same $\tau$.

\paragraph{Conclusion.}
The cycle-index and orbit formulas are two descriptions of the same invariant
dimension, and the explicit matrix calculations agree with both in every
tested case.

\end{document}
